% =====================================================================
\section{Ecuaciones de Ebers--Moll}
\label{sec:ebers-moll}
% =====================================================================

Sumando las componentes de la sección~\ref{sec:formal}, las corrientes de terminal quedan
\diapo{1.20}:
\[
  I_E = qA\left(\frac{D_E}{L_E}n_{E0}+\frac{D_B}{L_B}\frac{p_{B0}}{\tanh(W/L_B)}\right)\eVEB
      - qA\frac{D_B}{L_B}\frac{p_{B0}}{\sinh(W/L_B)}\eVCB
\]
\[
  I_C = qA\frac{D_B}{L_B}\frac{p_{B0}}{\sinh(W/L_B)}\eVEB
      - qA\left(\frac{D_C}{L_C}n_{C0}+\frac{D_B}{L_B}\frac{p_{B0}}{\tanh(W/L_B)}\right)\eVCB
\]
Tienen siempre la misma estructura: una combinación lineal de $\eVEB$ y $\eVCB$. Ponerles nombre
a los coeficientes da las ecuaciones de Ebers--Moll, válidas en \textbf{las cuatro regiones}.

\begin{clave}[Ecuaciones de Ebers--Moll \enhoja]
\[
  I_E = I_{F0}\eVEB - \aR I_{R0}\eVCB
\]
\[
  I_C = \aF I_{F0}\eVEB - I_{R0}\eVCB
\]
\[
  I_B = I_E - I_C
\]
con
\[
  I_{F0} = qA\left(\frac{D_E}{L_E}n_{E0}+\frac{D_B}{L_B}\frac{p_{B0}}{\tanh(W/L_B)}\right),
  \qquad
  I_{R0} = qA\left(\frac{D_C}{L_C}n_{C0}+\frac{D_B}{L_B}\frac{p_{B0}}{\tanh(W/L_B)}\right),
\]
\[
  \aF I_{F0} = \aR I_{R0} = qA\frac{D_B}{L_B}\frac{p_{B0}}{\sinh(W/L_B)}
  \qquad\text{(relación de reciprocidad)}.
\]
\end{clave}

\paragraph{Significado de cada parámetro.}
\begin{itemize}
  \item $I_{F0}$ ($=I_{ES}$ en otros libros): corriente de saturación de la juntura E--B
  medida con el colector en cortocircuito con la base ($V_{CB}=0$). Con $V_{CB}=0$, $I_E$ es la
  ecuación de un diodo ideal \diapo{1.20}.
  \item $I_{R0}$ ($=I_{CS}$): ídem para la juntura C--B con $V_{EB}=0$.
  \item $\aF$: ganancia de corriente en base común en modo directo ($\aF=I_C/I_E$ con
  $V_{CB}=0$). En MAD coincide con $\adc=\gamma\aT$ \diapo{1.22}.
  \item $\aR$: ídem en modo inverso (intercambiando emisor y colector). Es mucho menor
  que $\aF$ porque el colector está poco dopado (pésima ``eficiencia de emisión inversa'').
\end{itemize}

\begin{figure}[H]
\centering
\begin{circuitikz}[scale=0.95,transform shape]
  \draw (0,0) node[left] {E} to[short,o-] (1,0)
        to[D*,l_=$I_F$,-*] (3,0)
        to[D*,invert,l_=$I_R$,*-] (5,0) to[short,-o] (6,0) node[right] {C};
  \draw (1,0) -- (1,1.6) to[cI,l=$\aR I_R$,invert] (3,1.6) -- (3,0);
  \draw (3,0) -- (3,1.6) to[cI,l=$\aF I_F$] (5,1.6) -- (5,0);
  \draw (3,0) to[short,-o] (3,-1.3) node[below] {B};
  \draw[flecha] (0.2,-0.35) -- node[below,font=\small] {$I_E$} (0.8,-0.35);
  \draw[flecha] (5.2,-0.35) -- node[below,font=\small] {$I_C$} (5.8,-0.35);
  \draw[flecha] (2.7,-0.6) -- node[left,font=\small] {$I_B$} (2.7,-1.1);
\end{circuitikz}
\caption{Circuito equivalente de Ebers--Moll de un PNP (Pierret, fig.~11.6; \diapo{1.19, 1.22}).
$I_F=I_{F0}\eVEB$ e $I_R=I_{R0}\eVCB$ son las corrientes de dos diodos; cada diodo ``alimenta''
una fuente de corriente controlada sobre el otro terminal. El acoplamiento entre los dos diodos
es lo que los diferencia de ``dos diodos enfrentados''.}
\label{fig:em-circuito}
\end{figure}

\begin{figure}[H]
\centering
\begin{circuitikz}[scale=0.95,transform shape]
  \draw (0,0) node[left] {E} to[short,o-] (1,0)
        to[D*,l_=$I_F$,-*] (3,0)
        to[short] (5,0) to[short,-o] (6,0) node[right] {C};
  \draw (3,0) -- (3,1.6) to[cI,l=$\aF I_F$] (5,1.6) -- (5,0);
  \draw (3,0) to[short,-o] (3,-1.3) node[below] {B};
  \node[font=\small,align=left,anchor=west] at (6.8,0.4) {MAD: $I_R\approx-I_{R0}\approx0$\\$I_E\approx I_{F0}e^{V_{EB}/V_T}$\\$I_C\approx\aF I_E$};
\end{circuitikz}
\caption{En activa directa el diodo C--B (en inversa) y la fuente $\aR I_R$ son despreciables:
el transistor es \textbf{un diodo (E--B) más una fuente de corriente controlada} \diapo{1.23}.}
\label{fig:em-mad}
\end{figure}

\begin{ejemplo}[parámetros de Ebers--Moll]
Para el transistor de la sección~\ref{sec:parametros} ($L_C=\sqrt{D_C\tau_C}=\SI{59}{\micro m}$ con
$\tau_C=\SI{1}{\micro s}$):
\[
  I_{F0}=\SI{1.99e-14}{A},\quad I_{R0}=\SI{2.94e-14}{A},\quad
  \aF=0.9981,\quad \aR=\frac{\aF I_{F0}}{I_{R0}}=0.68 .
\]
En modo inverso $\beta_R=\aR/(1-\aR)\approx2$, frente a $\beta_F\approx525$ en modo directo:
el transistor es muy asimétrico.
\end{ejemplo}

\subsection{PNP y NPN}

Todo lo anterior vale para un NPN cambiando huecos por electrones e invirtiendo los signos
de todas las tensiones y corrientes.

\begin{table}[H]
\centering\small
\caption{Traducción PNP $\leftrightarrow$ NPN. Las diapositivas de la 2.\textsuperscript{a} parte
(tomadas de Sedra--Smith) usan NPN.}
\label{tab:pnp-npn}
\begin{tabular}{lll}
\toprule
 & PNP (Pierret, hoja) & NPN (Sedra) \\
\midrule
Portadores útiles en la base & huecos ($p_{B0}=n_i^2/N_B$) & electrones ($n_{B0}=n_i^2/N_A$) \\
Tensión de control en MAD & $V_{EB}>0$ & $V_{BE}>0$ \\
Juntura C--B en MAD & $V_{CB}<0$ & $V_{BC}<0$ \\
Tensión de salida & $V_{EC}>0$ & $V_{CE}>0$ \\
Corrientes ($>0$) & $I_E$ entra; $I_C$, $I_B$ salen & $I_C$, $I_B$ entran; $I_E$ sale \\
Ley exponencial & $I_C=I_Se^{V_{EB}/V_T}$ & $i_C=I_Se^{v_{BE}/V_T}$ \\
\bottomrule
\end{tabular}
\end{table}
